% LaTeX Template for Project Report, Version 2.0
% (Abstracted from a Major Project Report at CSED, NIT Calicut but can be
% modified easily to use for other reports also.)
%
% Released under Creative Commons Attribution license (CC-BY)
% Info: http://creativecommons.org/licenses/by/3.0/
%
% Created by: Kartik Singhal
% BTech CSE Batch of 2009-13
% NIT Calicut
% Contact Info: kartiksinghal@gmail.com
%
% It is advisable to learn the basics of LaTeX before using this template.
% A good resource to start with is http://en.wikibooks.org/wiki/LaTeX/
%
% All template fields are marked with a pair of angular brackets e.g. <title here>
% except for the ones defining citation names in ref.tex.
%
% Empty space after chapter/section/subsection titles can be used to insert text.
%
% Just compile this file using pdflatex after making all required changes.

\documentclass[12pt,a4paper]{report}
\usepackage[pdftex]{graphicx} %for embedding images
\usepackage{url} %for proper url entries
\usepackage[bookmarks, colorlinks=false, pdfborder={0 0 0}, pdftitle={Short and Long Term Goals}, pdfauthor={Athan L Clark}, pdfsubject={Report}, pdfkeywords={USMC}]{hyperref} %for creating links in the pdf version and other additional pdf attributes, no effect on the printed document
%\usepackage[final]{pdfpages} %for embedding another pdf, remove if not required
\usepackage{titlesec}
\usepackage{enumitem,amssymb}
\usepackage[normalem]{ulem}

\newlist{todolist}{itemize}{2}
\setlist[todolist]{label=$\square$}

\titleformat{\chapter}[hang]
  {\Large\bfseries}{\thechapter}{0.5em}{}

\begin{document}
\renewcommand\bibname{References} %Renames "Bibliography" to "References" on ref page

%include other pages
\begin{titlepage}

\begin{center}

\textup{\small {\bf Commandant's Reading List} \\ Book Report}\\[0.2in]

% Title
\Large \textbf {<Book Title here>}\\[0.5in]

       %% \small \emph{Submitted in partial fulfillment of\\
       %%  the requirements for the award of the degree of}
       %%  \vspace{.2in}

       %% {\bf Bachelor of Technology \\in\\ Computer Science and Engineering}\\[0.5in]

% Submitted by
\normalsize Submitted by \\
Pvt Clark, Athan L.

%% \vspace{.1in}
%% Under the guidance of\\
%% {\textbf{<Guide's name here>}}\\[0.2in]

\vfill

% Bottom of the page
\includegraphics[width=0.18\textwidth]{mcces-logo}\\[0.1in]
\Large{Marine Corps Communication Electronics School}\\
\normalsize
\textsc{Twentynine Palms}\\
California, 92278\\
\vspace{0.2cm}
\today{}

\end{center}

\end{titlepage}


\pagenumbering{roman} %numbering before main content starts
\tableofcontents

\newpage
\pagenumbering{arabic} %reset numbering to normal for the main content

\chapter{Summary}

<Summary here>
 %objective changed to problem definition
\chapter{Outline}

\section{Context}
\subsection{<any sub section here>}

\subsection{<related material>}

\subsubsection{<Sub-subsection title>}
some text\cite{citation-1-name-here}, some more text

\subsubsection{<Sub-subsection title>}
even more text\footnote{<footnote here>}, and even more.

\section{Setting}

\section{Characters}
\subsection{Lieutenant General Lord Wellington}

\subsection{General Craufurd}

\subsection{Rifleman Matthew Dodd}


\section{Events}

\section{Consequences}
 %literature survey included in this
\chapter{Analysis}

Here, I'll divulge into specific goals after digressing on the value of the statements
made in the previous chapter.

\section{USMC Related}

There are some very specific goals I'd like to achieve in my first (and possibly, only)
enlistment --- some of them for glory, and some cathartic; but overall, I would like to
see myself successful in this organization. A lot of it stems from the pride of my
family, but a lot comes from myself as well; joining the USMC feels like the best thing
I've ever done for myself (on good days\ldots on bad, it feels like the worst).

At any rate, being proud of my uniform and my belonging are fundamental
to my presence in the Corps, and I'll need to keep those values strong.


\begin{itemize}
\item Reach the rank of Sergeant --- My cousin is a Sergeant, but that's not why I want
  to be one\ldots I feel like its your \textit{responsibility} as enlisted to reach that rank;
  to be a \textit{servant} to the Corps (and a valued one, at that).
\item Achieve an associates degree in computer science and electrical engineering --- I already
  know so much about these fields, that they should be pretty easy to get. But, they would both
  be very valuable to my unit, and to myself. Likewise, they can be used in promotion boards,
  and generally speaking, would set me apart from many others.
\item Get jump wings --- this might not be possible for me, due to physical (and financial)
  concerns, but damn it would be cool to wear those on my cammies.
\item Earn a volunteer service medal --- Although not the most glorious medal to own, 100
  hours of volunteer service isn't a lot per-se, but it would make me feel good about myself.
\item I want to be well known in my shop and unit --- People should recognize me as someone
  trustworthy, hard working, strong, and admirable. I want people to respect me for where I've
  come from, and treat me like someone who earned it.
\item I should have 3 book reports done every year --- I'm already a year behind, so I better
  play catch-up!
\item I should take English 1 \& 2 --- GySgt Joiner gave us really good advice at the end
  of MAEMC\ldots I shouldn't pass that by.
\item I should have a ledger of impressionable material --- not specifically awards or letters
  of appreciation, but a record of my work done correctly, and more than necessary; the material
  that distinguishes me from others in my shop, solely from my craft.
\item Help others reach my rank --- it's my responsibility to pave the path for the next marines
  coming through my path; I have to give them a way to do so.
\item Buy a car --- it will be hard getting around without a vehicle, and owning one would
  get me a lot of freedom.
\item Bolster your TSP --- a retirement fund will be necessary 15 years from now; I would be
  foolish to waste that opportunity.
\item Save enough money to fund a company --- this will be very difficult, but a lifelong
  goal of mine; to own a successful company. It will be very hard to do that without finance,
  and a great way to get started is with the funding received from the military. Likewise,
  having the connections here may be essential to getting Local Cooking off the ground.
  \item Earn some cool challenge coins --- it's hard to describe the value of these.
\end{itemize}


\section{Interpersonal}

No man is an island; no one person can achieve something great on their own, and everyone has
had some help at some point in their journey. Working with other people is essential to succeeding in
any endeavor; military or otherwise. In order to achieve my long-term goal of starting a business,
I'll have to master the art of working with other people, but before that, I'll need to inspire my
Marines, and my chain of command's trust.

\begin{itemize}
\end{itemize}









\chapter{Conclusion}

<Conclusion here>

\cleardoublepage
%\pagebreak
\phantomsection
\addcontentsline{toc}{chapter}{References}
\begin{thebibliography}{99}

\bibitem{citation-1-name-here}<Name of the reference here>,\ \url{<url here>}

\bibitem{citation-2-name-here}<Name of the reference here>,\ \url{<url here>}

\end{thebibliography}

\chapter{Weekly Goals}

\section{20200420}

\begin{todolist}
\item Finish this document
\item Meet everyone in the shop
\item Get your bearings for where everything is located
\item Learn what the current state of affairs is in the shop
  \begin{todolist}
  \item What's busted?
  \item What's getting fixed?
  \item GCSS routine?
  \end{todolist}
\item Learn the routine
\item Find out expectations, and adhere to them
\item Find a way to be productive
\item Bring writing material
\end{todolist}

\chapter{Monthly Goals}

\section{April 2020}

\begin{todolist}
\item Get acquainted with Camp Legeune
\item Get some good workouts in
\item Restock on hygiene and cleaning supplies
\item Save some money
\item Pay off Puget Sound completely
\item Start payments for Cliffs Apartments
\item Learn liberty policies etc.
\item Get 5 sessions at the gym
\end{todolist}

\chapter{Quarterly Goals}

\section{April --- June 2020}

\begin{todolist}
\item Be goal and plan oriented, adjust your behavior
\item Focus on being limber and stretching
\item Meditate, and find a religious \& spiritual outlet
\item Write your letters
\item Treat your dwelling like a palace
\item Get a Navy Federal account, and plan on getting a credit card
\item Find out about your Legal Zoom payments, and LocalCooking LLC --- kill it, it's not needed
\item See the dentist, and the doctor for a checkup
\item Square away your uniforms --- get your charlies tailored, and a new barracks cover
\end{todolist}

\chapter{Yearly Goals}

\section{2020}

\begin{todolist}
\item Find some good friends, and maybe a girlfriend
\item Go to the beach, see what's around
\item Get over 280 on a PFT
\item Buy some good boots
\item Get your programming workstation in-order
\item Get some books from home
\item Get some more CLEP exams passed --- talk to an education center
\item Finish Leading Marines and the LCpl Seminar
\item Read 3 books on the Commandant's reading list, and write 3 reports
\item Volunteer at least 20 hours
\end{todolist}


\end{document}
